\section{The graph and its invariant subspaces}
\label{sec:graph}

Write $n=\prod_{i=1}^t p_i^{k_i}$, where the primes are distinct and
$k_i\geq1$. Denote by $G_n$ the graph whose vertices are the nonzero
proper ideals of $\mathbb Z_n$, with distinct ideals adjacent if their
intersection is nonzero. We use the combinatorial Laplacian $L=D-A$;
the graph is \emph{Laplacian integral} when every eigenvalue of $L$ is an
integer, including when the graph is disconnected.

The ideal $(\prod_i p_i^{\alpha_i})$ is encoded by
$0\leq\alpha_i\leq k_i$. We exclude $(0,\ldots,0)$, which encodes the
whole ring, and $(k_1,\ldots,k_t)$, which encodes the zero ideal.
Ideal intersection is generated by the least common multiple. Therefore
two distinct vertices $\alpha,\beta$ are adjacent exactly when
\[
\max(\alpha_i,\beta_i)<k_i\quad\hbox{for some }i.
\]
Equivalently, putting $S(\alpha)=\{i:\alpha_i<k_i\}$, adjacency is
$S(\alpha)\cap S(\beta)\ne\varnothing$.

For nonempty $S\subseteq[t]$, the support class $\mathcal C_S$ is a
clique of size $\prod_{i\in S}k_i$, except that the full support has
size $u=\prod_i k_i-1$. Edges between distinct support classes are complete
or absent according as their supports intersect or are disjoint. The full
support consists of universal vertices. Remove them and call the resulting
graph $H$; thus $G_n=K_u\vee H$, with the evident interpretation for $u=0$.

\begin{lemma}[Lifting across the universal vertices]
\label{lem:join}
If $L(H)v=\lambda v$ and $\lambda\ne0$, extension of $v$ by zero on
the $u$ universal vertices is an eigenvector of $L(G_n)$ with eigenvalue
$\lambda+u$.
\end{lemma}
\begin{proof}
Since $L(H)$ is symmetric with zero row sums, $\lambda\ne0$ implies
$\sum v=0$. At a vertex of $H$, its degree increases by $u$ and all
new neighbor values are zero. At a universal vertex the Laplacian action
is $-\sum v=0$. These are exactly the claimed eigenvector equations.
\end{proof}

For $(k_1,k_2,k_3)=(a,a,b)$, $H$ has six support classes of sizes
$a,a,b,a^2,ab,ab$ on $1,2,3,12,13,23$, respectively. The partition into
\[
X=\mathcal C_1\cup\mathcal C_2,\quad Y=\mathcal C_3,\quad
Z=\mathcal C_{12},\quad W=\mathcal C_{13}\cup\mathcal C_{23}
\]
is equitable, and on functions constant on these cells the Laplacian acts by
\begin{equation}
\label{eq:symmetric}
B=\begin{pmatrix}
a^2+ab&0&-a^2&-ab\\
0&2ab&0&-2ab\\
-2a&0&2a+2ab&-2ab\\
-a&-b&-a^2&a^2+a+b
\end{pmatrix}.
\end{equation}
For example, an $X$ vertex has $a^2$ neighbors in $Z$ and $ab$ in $W$,
while a $W$ vertex has $a,b,a^2$ neighbors in $X,Y,Z$.
Internal neighbors cancel for cell-constant functions. The cell sizes are
$2a,b,a^2,2ab$, and multiplying $B$ on the left by their diagonal matrix
gives a symmetric matrix. Thus $B$ is similar to a real symmetric matrix.

The functions taking values $s,-s$ on $\mathcal C_1,\mathcal C_2$,
$t,-t$ on $\mathcal C_{13},\mathcal C_{23}$, and zero elsewhere
form a second invariant subspace, with matrix
\begin{equation}
\label{eq:antisymmetric}
B_-=\begin{pmatrix}a^2+ab&-ab\\-a&a^2+a+b+2ab\end{pmatrix}.
\end{equation}
These functions have sum zero. Their eigenvalues therefore also lift by
$u=a^2b-1$, directly by the proof of Lemma~\ref{lem:join}.

Expanding~\eqref{eq:symmetric} gives $\det(xI-B)=x f_{a,b}(x)$, where
\begin{align}
f_{a,b}(x)&=x^3-E_1x^2+E_2x-E_3,\label{eq:cubic}\\
E_1&=2a^2+5ab+3a+b,\notag\\
E_2&=a^4+7a^3b+3a^3+8a^2b^2+11a^2b+2a^2+3ab^2+2ab,\notag\\
E_3&=2a^2b(a+b+1)(a^2+2ab+2a+b).\notag
\end{align}
Its roots are real and nonzero: reality follows from the similarity above,
and $E_3>0$. Every root $\lambda$ gives an eigenvalue $\lambda+u$ of
$G_n$. A noninteger root of $f_{a,b}$ therefore proves nonintegrality.
\section{Repeated exponents}
\label{sec:bounds}
\label{sec:completion}
Coordinate interchange gives the two invariant blocks above. The
antisymmetric discriminant handles most parameters; exact factor indices
and the symmetric cubic settle every square-discriminant exception.

\subsection{Two unit exponents}

\begin{theorem}
\label{thm:pqrk}
For distinct primes $p,q,r$ and every $k\geq1$, the graph $G_{pq r^k}$
has the irrational eigenvalues
\[
\frac{6k+1\pm\sqrt{(2k+2)^2-3}}2.
\]
In particular, every permutation of the exponent vector $(1,1,k)$
gives a nonintegral Laplacian spectrum.
\end{theorem}
\begin{proof}
The support classes $1,2$ have size $1$, and $13,23$ have size $k$.
Consider functions with values $s,-s$ on the first two classes,
$t,-t$ on the latter two, and zero elsewhere.
A vertex of support $1$ has $2k$ neighbors and Laplacian value
$2ks-kt$. A vertex of support $13$ has degree $4k$; after subtracting
the same-class and opposite-class values its Laplacian value is
$-s+(4k+1)t$. The swapped classes give their negatives, and all other
classes give zero by cancellation. The restricted matrix of the full
graph is therefore
\[
B_k=\begin{pmatrix}2k&-k\\-1&4k+1\end{pmatrix}.
\]
Its characteristic polynomial is $x^2-(6k+1)x+8k^2+k$, with
discriminant $(2k+2)^2-3$. Since
\[
(2k+1)^2<(2k+2)^2-3<(2k+2)^2
\]
for $k\geq1$, the two roots are irrational. Relabeling prime coordinates
preserves the graph.
\end{proof}
\subsection{The boundary family}
The following argument was supplied by Reza Nikandish in the public
collaboration discussion. It is the index-zero case in the factor
reduction below.

\begin{theorem}
\label{thm:boundary}
For every $a\geq2$, put $C=(a-1)(2a-1)$. For distinct primes $p,q,r$,
the graph $G_{p^a q^a r^C}$ is not Laplacian integral, although the
antisymmetric block~\eqref{eq:antisymmetric} has integer eigenvalues.
\end{theorem}
\begin{proof}
Specialize $f_{a,b}$ from~\eqref{eq:cubic} to $b=C$.
For $a=2$ it is $x^3-47x^2+702x-3312$. Its values at
$0,1,2,3,4$ modulo $5$ are $3,4,2,3,3$. Thus it has no root in
$\mathbb F_5$ and is irreducible over $\mathbb Q$; its real roots
are noninteger. Lemma~\ref{lem:join} gives noninteger graph eigenvalues.

For $a\geq3$ put $m=2a^3-2a^2+2a-1$. Expansion gives
\begin{align*}
f_{a,C}(m)&=2(5a^4-10a^3+7a^2-1),\\
f_{a,C}(m-1)&=-2(a-1)(a+2)(2a^4-7a^3+7a^2-a-3).
\end{align*}
The first value is positive because its parenthesis is
$5a^3(a-2)+7a^2-1$. Writing $z=a-3\geq0$, the last factor in the
second value is
$2z^4+17z^3+52z^2+68z+30>0$.
Consequently $f_{a,C}(m-1)<0<f_{a,C}(m)$, so a real root lies strictly
between $m-1$ and $m$. Its integer shift by $u=a^2C-1$ is noninteger.

The characteristic discriminant of $B_-$ is
$D=(a+1)^2b^2+2a(3a+1)b+a^2$.
At $b=C$ it equals $(2a^3-a^2+a-1)^2$, and the two eigenvalues are
$4a^3-3a^2+a$ and $2a^3-2a^2+1$. Thus the obstruction in this family
comes from the symmetric cubic.
\end{proof}
\subsection{Exact factor pairs and the first indices}

Fix $a,b\geq1$, and put
\[
h=a+1,\quad C=(a-1)(2a-1),\quad M=a^3(2a+1),\quad
A=h^2b+a(3a+1).
\]
The characteristic discriminant of~\eqref{eq:antisymmetric} is
\begin{equation}
\label{eq:discriminant}
D=h^2b^2+2a(3a+1)b+a^2.
\end{equation}
If $D$ is not a square, this block has irrational eigenvalues, which
remain irrational after the integer shift $a^2b-1$.

\begin{proposition}[Exact divisor reduction]
\label{prop:divisors}
For fixed $a\geq1$, the positive integers $b$ for which $D$ is a square
are exactly those obtained from positive factor pairs $d\leq e$ of $M$
such that
\begin{equation}
\label{eq:divisor-conditions}
b=\frac{d+e-a(3a+1)}{h^2}\in\mathbb Z_{\geq1},\qquad
s=\frac{e-d}{h}\in\mathbb Z_{\geq0}.
\end{equation}
In that case $D=s^2$.
\end{proposition}
\begin{proof}
The polynomial identity
\begin{equation}
\label{eq:product}
A^2-h^2D=4M
\end{equation}
holds for arbitrary $a,b$. If $D=s^2$ with $s\geq0$, then
$A-hs$ and $A+hs$ are positive: their product is positive and their sum
is $2A>0$. They have the same parity; since their product is divisible
by $4$, both are even. Write $A-hs=2d$ and $A+hs=2e$.
Then $de=M$, $d+e=A$ and $e-d=hs$, giving~\eqref{eq:divisor-conditions}.
Conversely those conditions yield $A=d+e$ and $hs=e-d$.
Using $(d+e)^2-(e-d)^2=4de=4M$ in~\eqref{eq:product} gives $D=s^2$.
\end{proof}

The divisor set is finite for every fixed $a$; this criterion has no
truncation in $b$. It gives first a quadratic bound.

\begin{corollary}
\label{cor:quadratic}
If $D$ is a square then $b\leq C$. Consequently every $b>C$ gives
a nonintegral Laplacian spectrum. For $a=1$, every $b\geq1$ is covered.
\end{corollary}
\begin{proof}
Since $(d-1)(e-1)\geq0$, $A=d+e\leq M+1$. The identity
\[
M+1-a(3a+1)=h^2C
\]
then gives $b\leq C$. The nonsquare case is covered by the antisymmetric
block. At $a=1$, $C=0$.
\end{proof}

\begin{lemma}
\label{lem:congruence}
Every admissible factor pair in Proposition~\ref{prop:divisors} satisfies
$d\equiv e\equiv1\pmod h$.
\end{lemma}
\begin{proof}
From $e-d=hs$ and $d+e=A$, reduction modulo $h$ gives
$e\equiv d$ and $2d\equiv2$, since $a(3a+1)\equiv2\pmod h$.
Also
\begin{equation}
\label{eq:strong-product}
(d-1)(e-1)=M-A+1=h^2(C-b).
\end{equation}
If $h$ is odd, the congruence $2(d-1)\equiv0$ gives the assertion.
If $h$ is even, the alternative to $d\equiv1$ is
$d-1\equiv h/2\pmod h$. Then $d-1$ and $e-1$ would both equal
$h/2$ times an odd integer. Their product divided by $h^2$ would be an
odd integer divided by $4$, contradicting~\eqref{eq:strong-product}.
This excludes the alternative for either parity of $a$.
\end{proof}

The first possible value after $d=1$ is $d=a+2$.
Since $d$ divides $M$, the identity
\[
M=(a+2)(2a^3-3a^2+6a-12)+24
\]
forces $a+2\mid24$. With $a\geq2$ this leaves
$a=2,4,6,10,22$. Substitution into~\eqref{eq:divisor-conditions}
gives $b=0,2,5,12,35$, respectively. Excluding $b=0$, the complete
positive-$b$ list and its cubic certificates are:
\begin{center}
\begin{tabular}{clc}
\toprule
$(a,b)$ & $f_{a,b}(x)$ & Root-free prime\\
\midrule
$(4,2)$ & $x^3-86x^2+2304x-18816$ & 5\\
$(6,5)$ & $x^3-245x^2+19266x-488160$ & 7\\
$(10,12)$ & $x^3-842x^2+230160x-20534400$ & 13\\
$(22,35)$ & $x^3-4919x^2+7887858x-4132479120$ & 17\\
\bottomrule
\end{tabular}
\end{center}
Each polynomial has no root modulo its indicated prime. Its monic cubic
is therefore irreducible over $\mathbb Q$, and its real nonzero roots
lift to noninteger graph eigenvalues. Thus every possible $d=a+2$ pair
is settled, over all $a$, by this finite divisor argument.

\subsection{A finite reduction for each factor index}

For a nonboundary square-discriminant pair, write $d=1+hj$, $j\geq1$.
The following criterion fixes $j$ instead of $a$.

\begin{proposition}
\label{prop:index}
Fix $j\geq1$ and set $K_j=(j+1)^3(j+2)$.
Every admissible square-discriminant pair of index $j$ is obtained from
a positive divisor $d$ of $K_j$ with
$a=(d-1)/j-1\in\mathbb Z_{\geq2}$.
Set $e=M/d$, $v=(e-1)/(a+1)$, and retain exactly those candidates with
$d\leq e$ and $b=C-jv\geq1$. Then $e,v,b$ are integers and
$D=(v-j)^2$. Conversely every retained candidate is admissible.
\end{proposition}
\begin{proof}
Modulo $d=j(a+1)+1$ we have $ja\equiv-(j+1)$, whence
\[
j^4M=(ja)^3j(2a+1)\equiv(j+1)^3(j+2)=K_j\pmod d.
\]
Since $\gcd(d,j)=1$, we have $d\mid M$ if and only if $d\mid K_j$.
Thus $e$ is integral. Since $M\equiv d\equiv1\pmod h$, also
$e\equiv1\pmod h$, so $v$ is integral. The identity
\[
M=1+h(3a-2)+h^2C=(1+hj)(1+hv)
\]
gives $j+v+h jv=3a-2+hC$. Defining $b=C-jv$, we obtain
$d+e=h^2b+a(3a+1)$ and $e-d=h(v-j)$.
The ordering condition gives $v-j\geq0$; Proposition~\ref{prop:divisors}
now proves admissibility and $D=(v-j)^2$. The reverse direction follows
by applying these identities to any admissible index-$j$ pair.
\end{proof}

This is a complete finite problem for each fixed $j$, without a
truncation in $a$. Section~\ref{sec:completion} combines the first
three such problems with a uniform argument for all later indices.
For $j=2$, $K_2=108$ and $d=2a+3$ is an odd divisor of $108$ with
$d\geq7$. The only candidates are $d=9,27$, giving $a=3,12$.
The first has $e=21,v=5,b=0$ and is excluded. The second has
$e=1600,v=123,b=7$ and $\sqrt D=121$.
Thus $(12,7)$ is the only positive-$b$ case with $j=2$.

Its symmetric cubic is $f_{12,7}(x)=x^3-751x^2+180348x-13829760$.
Its values at all residues modulo $13$ are
$4,7,9,3,8,4,10,6,11,5,7,10,7$.
As a monic cubic with no root modulo $13$, it is irreducible over
$\mathbb Q$. Its real nonzero roots lift to noninteger graph eigenvalues,
settling the entire index-$2$ case.

\subsection{Completion by a cubic interval}

\begin{theorem}
\label{thm:repeated}
For distinct primes $p,q,r$ and all integers $a,b\geq1$, the graph
$G_{p^a q^a r^b}$ is not Laplacian integral. The same conclusion holds
for every permutation of the exponent vector $(a,a,b)$.
\end{theorem}

The arithmetic reduction and a uniform sign obstruction for the cubic
combine to exhaust the family. We first settle one last finite factor level.

\begin{lemma}
\label{lem:third-index}
Every square-discriminant pair with index $j=3$ gives a nonintegral
Laplacian spectrum.
\end{lemma}
\begin{proof}
Proposition~\ref{prop:index} gives $d\mid K_3=320$ and $d=3a+4$.
For $a\geq2$, the complete divisor list of this shape is
$10,16,40,64,160$. The resulting values are
\begin{center}
\begin{tabular}{rrrrrl}
\toprule
$d$ & $a$ & $e$ & $v$ & $b$ & Outcome\\
\midrule
10 & 2 & 4 & 1 & 0 & $d>e$\\
16 & 4 & 36 & 7 & 0 & $b=0$\\
40 & 12 & 1080 & 83 & 4 & Root-free modulo $5$\\
64 & 20 & 5125 & 244 & 9 & Root-free modulo $7$\\
160 & 52 & 92274 & 1741 & 30 & Root-free modulo $7$\\
\bottomrule
\end{tabular}
\end{center}
The three positive-$b$ cubics are
\begin{align*}
f_{12,4}(x)&=x^3-568x^2+100032x-5248512,\\
f_{20,9}(x)&=x^3-1769x^2+992820x-174744000,\\
f_{52,30}(x)&=x^3-13394x^2+57771168x-80229951360.
\end{align*}
Their values at all residues of the indicated primes are, respectively,
$(3,3,3,4,2)$, $(3,2,4,1,6,4,1)$ and $(1,6,4,1,3,2,4)$.
Each monic cubic is therefore irreducible over $\mathbb Q$ and has
noninteger real roots. Section~\ref{sec:graph} supplies the lift to
noninteger graph eigenvalues.
\end{proof}

\begin{lemma}[A uniform consecutive-integer interval]
\label{lem:sign}
For $a,b\geq1$, put
\[
L(a)=\frac{a^3+2a^2+2a+1}{a(2a+1)}.
\]
If $b<L(a)$, then $f_{a,b}$ has a root strictly between $2ab-1$ and $2ab$.
\end{lemma}
\begin{proof}
Substitution in~\eqref{eq:cubic} gives
\begin{align*}
f_{a,b}(2ab)&=2a^2b^2>0,\\
f_{a,b}(2ab-1)&=-(a+b+1)
\bigl(a^3+2a^2+2a+1-a(2a+1)b\bigr).
\end{align*}
The second value is strictly negative when $b<L(a)$. The intermediate
value theorem gives the asserted noninteger, nonzero root.
\end{proof}

\begin{proof}[Proof of Theorem~\ref{thm:repeated}]
For $a=1$, Corollary~\ref{cor:quadratic} settles every positive $b$.
For $a\geq2$, nonsquare $D$ gives irrational antisymmetric eigenvalues.
When $D$ is square, index $j=0$ is Theorem~\ref{thm:boundary},
and indices $j=1,2$ are settled in Section~\ref{sec:bounds}.
Lemma~\ref{lem:third-index} settles $j=3$.

It remains to consider $j\geq4$. Then $d\geq Q=4a+5$, and
$e\geq d\geq Q$ implies
\[
\frac MQ+Q-(d+e)=(d-Q)\left(\frac eQ-1\right)\geq0.
\]
Since $d+e=(a+1)^2b+a(3a+1)$, we obtain
\[
b\leq B_4(a)=\frac{M/Q+Q-a(3a+1)}{(a+1)^2}
=\frac{(a-5)(2a-5)}{4a+5}.
\]
For every $a\geq2$,
\[
L(a)-B_4(a)=
\frac{41a^3-17a^2-11a+5}{a(2a+1)(4a+5)}>0.
\]
Indeed, with $z=a-2\geq0$, the numerator becomes
$41z^3+229z^2+413z+243$. Thus $b\leq B_4(a)<L(a)$, and
Lemma~\ref{lem:sign} gives a noninteger cubic root, which lifts across
the universal vertices. These cases exhaust all positive $a,b$.
Permuting the prime factors preserves the graph's isomorphism type.
\end{proof}

\section{The complement quotient and small minima}
\label{sec:unit}
For unequal exponents the complement quotient retains all six support
classes. Its positive roots reflect to graph eigenvalues by the complement
identity and the universal-vertex lift.

\subsection{A unit exponent}

\begin{theorem}
\label{thm:unit}
For distinct primes $p,q,r$ and every $b,c\geq1$, the graph
$G_{p q^b r^c}$ is not Laplacian integral. The same holds for
permutations of the exponent vector $(1,b,c)$.
\end{theorem}

The equal-exponent decomposition is replaced by a complement identity
on all six support classes. For arbitrary $a,b,c\geq1$, remove the
$u=abc-1$ universal vertices and denote the remaining graph by $H$.
In the order $1,2,3,12,13,23$, its class-size vector is
$w=(a,b,c,ab,ac,bc)^T$ and its vertex count is
$N=a+b+c+ab+ac+bc$. Disjoint supports are adjacent in the complement
$\overline H$, whose cell-constant Laplacian quotient is
\[
C=\begin{pmatrix}
b+c+bc&-b&-c&0&0&-bc\\
-a&a+c+ac&-c&0&-ac&0\\
-a&-b&a+b+ab&-ab&0&0\\
0&0&-c&c&0&0\\
0&-b&0&0&b&0\\
-a&0&0&0&0&a
\end{pmatrix}.
\]
Multiplication on the left by $\operatorname{diag}(w)$ makes $C$
symmetric. Moreover $C\mathbf1=0$ and $w^TC=0$.
If $B$ is the quotient of $H$, the complement relation gives
\begin{equation}
\label{eq:complement}
B+C=NI-\mathbf1w^T.
\end{equation}
Thus an eigenvector $Cv=\mu v$ with $\mu\ne0$ satisfies $w^Tv=0$
and $Bv=(N-\mu)v$. When $N-\mu\ne0$, Lemma~\ref{lem:join}
lifts it to an eigenvalue $N+u-\mu$ of $G_n$.

\begin{proof}[Proof of Theorem~\ref{thm:unit}]
Write $\det(xI-C)=x h_{a,b,c}(x)$. Direct determinant expansion yields
\begin{align*}
h_{a,b,c}(0)&=-abc(a+b+c)N,\\
h_{a,b,c}(a)&=-a^2b^2c^2(a^2+2a-b-c).
\end{align*}
At $a=1$, the first value is strictly negative and
\[
h_{1,b,c}(1)=b^2c^2(b+c-3).
\]
When $b+c>3$, a real root $\mu$ lies strictly between $0$ and $1$.
Equation~\eqref{eq:complement} and the universal-vertex lift produce a
Laplacian eigenvalue $V-\mu\in(V-1,V)$, where $V=N+u$ is the integer
vertex count. This eigenvalue is noninteger. When $b+c\leq3$, the
positive pairs are $(1,1),(1,2),(2,1)$; Theorem~\ref{thm:pqrk}
settles the corresponding permutations of $(1,1,k)$. Permuting the
prime factors preserves the graph's isomorphism type.
\end{proof}

\subsection{A minimum exponent of two}\label{sec:minimum-two}

\begin{theorem}
\label{thm:minimum-two}
For distinct primes $p,q,r$ and all $b,c\geq1$, the graph
$G_{p^2q^br^c}$ is not Laplacian integral. Consequently every
three-prime exponent vector with minimum entry at most two gives a
nonintegral graph.
\end{theorem}

\begin{proof}
Use the complement quotient and quintic from Section~\ref{sec:unit}.
For $a=2$, $h_{2,b,c}(0)<0$. Putting $s=b+c$ and $t=bc$ gives
\[
h_{2,b,c}(1)=(s-13)t^2-(2s+6)t+3s^3+s^2-3s-1.
\]
By symmetry order $b\leq c$. A unit exponent is covered by
Theorem~\ref{thm:unit}; if $b=2$ or $b=c$, use
Theorem~\ref{thm:repeated}. It remains to consider $3\leq b<c$.

When $b=3$ and $c=5+v$, $v\geq0$, the exact expansion is
\[
h_{2,3,5+v}(1)=12v^3+112v^2+268v+120>0.
\]
When $b\geq4$ and $c>b$, write $b=4+u$ and $c=5+u+v$, with
$u,v\geq0$. Then
\begin{align*}
h_{2,4+u,5+u+v}(1)
={}&(u^2+8u+19)v^3\\
&+(4u^3+38u^2+128u+170)v^2\\
&+(5u^4+62u^3+271u^2+502u+368)v\\
&+2u^5+32u^4+190u^3+504u^2+552u+160>0.
\end{align*}
All coefficients are positive. In either region the opposite signs at
$0$ and $1$ give a root $\mu\in(0,1)$ and, by~\eqref{eq:complement}
and the universal-vertex lift, a noninteger graph eigenvalue in $(V-1,V)$.

The only remaining pair is $(b,c)=(3,4)$. Here
\[
h_{2,3,4}(x)=x^5-53x^4+953x^3-6523x^2+13134x-7560,
\]
and $h_{2,3,4}(1)=-48$ while
$h_{2,3,4}(3/2)=13437/32>0$. A root lies in $(1,3/2)$.
Since $V=58$, its lifted graph eigenvalue lies in $(113/2,57)$,
an interval containing no integer. This exhausts all positive $b,c$.
Combining with Theorem~\ref{thm:unit} proves the minimum-entry assertion.
\end{proof}
\subsection{The uniform cutoff and minimum three}\label{sec:distinct-tail}

\begin{proposition}
\label{prop:distinct-tail}
For distinct primes $p,q,r$ and positive exponents ordered as
$2\leq a\leq b\leq c$, the graph $G_{p^a q^b r^c}$ is not
Laplacian integral whenever $c\geq4a^2-2a$.
\end{proposition}

\begin{proof}
Use the complement quintic $h=h_{a,b,c}$ of Section~\ref{sec:unit}.
Its first endpoint is the cubic
\begin{equation}
\label{eq:general-first-endpoint}
h(1)=Lc^3+Qc^2+Rc+S,
\end{equation}
where
\begin{align*}
L={}&a^2+b^2-1,\\
Q={}&a^3-4a^2b^2+2a^2b-a^2+2ab^2+ab-3a
       +b^3-b^2-3b+3,\\
R={}&(a-1)(b-1)(2ab+3a+3b-3),\\
S={}&(a-1)(b-1)(a+b-1)(ab+a+b-1).
\end{align*}
Here $L,R,S>0$. In fact the sharper condition $Lc+Q\geq0$
already gives $h(1)>0$. Put $T(a)=4a^2-2a$. Then
\[
Q+T(a)L=b^2(b-1)+(2a^2+a-3)b+K(a),
\]
where $K(a)=4a^4-a^3-5a^2-a+3$. For $a=2+u$, $u\geq0$,
\[
K(2+u)=4u^4+31u^3+85u^2+95u+37>0.
\]
Consequently $Lc+Q>0$ for $c\geq T(a)$. Since $h(0)<0$,
a root lies in $(0,1)$, producing a noninteger graph eigenvalue in
$(V-1,V)$ by~\eqref{eq:complement} and Lemma~\ref{lem:join}.
\end{proof}

\begin{theorem}
\label{thm:minimum-three}
For distinct primes $p,q,r$ and all $b,c\geq1$, the graph
$G_{p^3q^br^c}$ is not Laplacian integral. Consequently every
three-prime exponent vector with minimum entry at most three gives a
nonintegral graph.
\end{theorem}

\begin{proof}
Theorems~\ref{thm:unit}, \ref{thm:minimum-two} and \ref{thm:repeated}
leave $4\leq b<c$. Proposition~\ref{prop:distinct-tail} settles
$c\geq30$. At $a=3$, equation~\eqref{eq:general-first-endpoint} is
\begin{align*}
h_{3,b,c}(1)={}&(b^2+8)c^3+(b-1)(b^2-30b-12)c^2\\
&+6(b-1)(3b+2)c+4(b-1)(b+2)(2b+1).
\end{align*}
The remaining domain $4\leq b<c<30$ has exactly
$\binom{26}{2}=325$ pairs. Exhaustive integer evaluation of this
cubic gives $257$ positive and $68$ negative values, with no zero.
Table~\ref{tab:minimum-three-signs} lists every negative range;
all other pairs in this finite domain have positive values.
The complete polynomial and separate Horner sign certificates are in
\path{results/distinct-tail.json}; their construction is described in
Section~\ref{sec:verification}. Thus a positive value gives a root in $(0,1)$.

\begin{table}[ht]
\centering
\begin{tabular}{rrrrrrrrrrr}
\toprule
$b$&4&5&6&7&8&9&10&11&12&13\\
\midrule
$c_{\min}$&5&6&7&8&9&10&11&12&13&14\\
$c_{\max}$&13&15&16&17&17&16&16&15&14&14\\
\bottomrule
\end{tabular}
\caption{Every negative $h_{3,b,c}(1)$ in $4\leq b<c<30$;
each column includes all integers between its two endpoints.}
\label{tab:minimum-three-signs}
\end{table}

For the negative values, the second endpoint is positive except at
$(b,c)=(4,5)$. Indeed, for $v\geq0$,
\[
h_{3,4,6+v}(2)=104v^3+1238v^2+4018v+2344>0.
\]
For $b\geq5$, write $b=5+u$, $c=6+u+v$, with $u,v\geq0$.
Direct expansion gives
\begin{align*}
h_{3,5+u,6+u+v}(2)
={}&(5u^2+54u+153)v^3\\
&+(20u^3+262u^2+1179u+1863)v^2\\
&+(25u^4+426u^3+2642u^2+7011u+6624)v\\
&+10u^5+218u^4+1828u^3+7200u^2+12600u+6480>0.
\end{align*}
These regions exhaust $4\leq b<c$ except $(4,5)$. Whenever
$h(1)<0<h(2)$, a root lies in $(1,2)$, giving a noninteger graph
eigenvalue in $(V-2,V-1)$.

Finally
\[
h_{3,4,5}(x)=x^5-83x^4+2327x^3-24133x^2+60576x-42480.
\]
Here $h(1)=-3792$ and $h(3/2)=48825/32>0$. The root in
$(1,3/2)$ lifts to a graph eigenvalue in $(233/2,117)$, since $V=118$.
Every interval used contains no integer. The finite nonvanishing
certificate and the written infinite-range arguments exhaust all $b,c$;
Theorems~\ref{thm:unit} and \ref{thm:minimum-two} give the final assertion.
\end{proof}
\section{Small-minimum certificates and a uniform low root}
\label{sec:low-spectrum}
The complement quotient also makes sense for positive real support weights.
With $W=\operatorname{diag}(w)$ and disjointness edges on the six supports,
\[
 z^TWCz=\sum_{\{S,T\}:\,S\cap T=\varnothing}w_Sw_T(z_S-z_T)^2.
\]
Thus $W^{1/2}CW^{-1/2}$ is positive semidefinite. The singleton triangle
and the three complementary-pair edges connect this support graph, so
its kernel is one-dimensional. This justifies the real-parameter spectral
arguments below as well as their integer specializations.

The uniform cutoff leaves a complete finite domain at minima four through
seven. A symmetric Schur complement then supplies a positive root below
three at every larger minimum, including when an endpoint sign change
misses a pair of low roots.

\subsection{The complete minimum-through-seven certificate}

The uniform cutoff makes the requested diagnostic range complete for each
of the minimum exponents $4,5,6,7$. Here the relevant quintic is
$g_B(x)=\det(xI-B)/x$ for $H$, whereas the preceding endpoint arguments
use the complement quintic $h_C(x)=-g_B(N-x)$.

\begin{corollary}
\label{cor:minimum-seven}
Every positive three-prime exponent triple with minimum entry at most
seven gives a nonintegral ideal intersection graph.
\end{corollary}

\begin{proof}
Theorem~\ref{thm:minimum-three} handles minimum entries at most three.
For a minimum entry $a\in\{4,5,6,7\}$, repeated entries are covered by
Theorem~\ref{thm:repeated}, and Proposition~\ref{prop:distinct-tail}
settles $c\geq4a^2-2a$. The remaining ordered triples lie in the exact
finite domain $a<b<c<4a^2-2a$.

For all $27\,562$ triples in this domain, the integer discriminant
$\Delta_{a,b,c}$ of $g_B$ is positive and strictly between two
consecutive integer squares. Table~\ref{tab:open-region} gives the
complete counts. The certificate
\path{results/open-region-discriminants.csv} records every triple,
discriminant and integer $q$ satisfying
$q^2<\Delta_{a,b,c}<(q+1)^2$. Each discriminant is independently
recomputed as a $9\times9$ integer Sylvester determinant by
\path{scripts/check_open_region_certificate.py}, using fraction-free
elimination and exact divisions. The certificate exhausts the domain
without omitted or duplicate triples.

If all roots of the monic quintic $g_B$ were integers, its discriminant,
the product of the squared pairwise root differences, would be an
integer square. Thus every certified quintic has a noninteger root.
The quotient is similar to a real symmetric Laplacian, so the root is
real and nonzero; Lemma~\ref{lem:join} gives a noninteger graph
eigenvalue. This complete finite certificate and the written cutoff
exhaust all triples with minimum entry at most seven.
\end{proof}

\begin{table}[ht]
\centering
\begin{tabular}{rrrrrr}
\toprule
Minimum $a$&4&5&6&7&Total\\
\midrule
Exclusive cutoff&56&90&132&182&---\\
Complete pairs&1275&3486&7750&15051&27562\\
Nonsquare discriminants&1275&3486&7750&15051&27562\\
\bottomrule
\end{tabular}
\caption{Exhaustive discriminant certificates in the region derived from
the uniform cutoff.}
\label{tab:open-region}
\end{table}

\subsection{The symmetric Schur complement}

\begin{proposition}
\label{prop:low-spectrum}
Order the exponents as $a\leq b\leq c$, and let $m$ be an integer
with $2\leq m<a$. If
\begin{equation}
\label{eq:upper-inertia}
(c-m)(a+b+c-m)<mab,
\end{equation}
then the six-support complement quotient has exactly two nonzero
eigenvalues in $(0,m)$. If, in addition,
\begin{equation}
\label{eq:lower-inertia}
(a-m+1)(a+b+c-m+1)>(m-1)bc,
\end{equation}
both lie in $(m-1,m)$, so $G_{p^a q^b r^c}$ is not Laplacian integral.
\end{proposition}

\begin{proof}
Put $s=a+b+c$ and $P=abc$. Symmetrize $C$ by the square roots of its
cell sizes and reorder the classes as $1,2,3,23,13,12$. The resulting
real symmetric matrix is
\[
\mathcal C=\begin{pmatrix}A&-\sqrt P I\\-\sqrt P I&D\end{pmatrix},
\quad D=\operatorname{diag}(a,b,c),
\]
where $A$ has diagonal $(b+c+bc,a+c+ac,a+b+ab)$ and
off-diagonal entries $A_{ij}=-\sqrt{a_i a_j}$, with
$(a_1,a_2,a_3)=(a,b,c)$. For $0<x<a$, the block $D-xI$ is positive
definite, and Sylvester's law of inertia gives
\[
\operatorname{inertia}(\mathcal C-xI)
=\operatorname{inertia}(D-xI)+\operatorname{inertia}(K(x)),
\]
where
\[
K(x)=\operatorname{diag}(k_a(x),k_b(x),k_c(x))-vv^T,
\quad v=(\sqrt a,\sqrt b,\sqrt c)^T,
\quad k_t(x)=s-x-\frac{xP}{t(t-x)}.
\]
For fixed $x>0$, the function $k_t(x)$ increases with $t>x$, since
\[
\frac{\partial k_t(x)}{\partial t}
=\frac{xP(2t-x)}{t^2(t-x)^2}>0.
\]
Condition~\eqref{eq:upper-inertia} is exactly $k_c(m)<0$.
Thus all three diagonal values are negative, so $K(m)$ is negative
definite. The matrix $\mathcal C-mI$ has three negative and three
positive eigenvalues. The complement support graph is a triangle with
one pendant cell at each vertex; it is connected, so $C$ has a simple
zero eigenvalue and all other eigenvalues are positive. Exactly two
nonzero eigenvalues therefore lie in $(0,m)$.

For $x=m-1$, condition~\eqref{eq:lower-inertia} says $k_a(x)>0$,
so all three $k_t(x)$ are positive. Congruence by the inverse square
roots of these values turns $K(x)$ into $I-zz^T$, where
\[
\|z\|^2=\frac{a}{k_a(x)}+\frac{b}{k_b(x)}+\frac{c}{k_c(x)}>1,
\]
because each $k_t(x)<s$. Consequently $K(x)$ has one negative and two
positive eigenvalues, and the only eigenvalue of $C$ below $x$ is zero.
Both low positive eigenvalues lie in $(m-1,m)$. They are less than
$a<N$, so the complement identity and universal-vertex lift apply.
Their graph eigenvalues lie in $(V-m,V-m+1)$, which contains no integer.
\end{proof}

\subsection{The uniform comparison at three}

\begin{theorem}
\label{thm:uniform-low-root}
For every ordered triple $4\leq a\leq b\leq c$, the six-support
complement quotient has a nonzero eigenvalue in $(0,3)$.
\end{theorem}

\begin{proof}
Use $s=a+b+c$, $P=abc$ and $v=(\sqrt a,\sqrt b,\sqrt c)^T$ from
Proposition~\ref{prop:low-spectrum}, and define
\[
L=\operatorname{diag}\left(s-\frac{3P}{a^2},
s-\frac{3P}{b^2},s-\frac{3P}{c^2}\right)-vv^T.
\]
Expansion of this diagonal-minus-rank-one determinant gives
\begin{align*}
\det L&=6\bigl(c(a-b)^2+b(a-c)^2+a(b-c)^2\bigr),\\
v^TLv&=-3P\left(\frac1a+\frac1b+\frac1c\right)<0.
\end{align*}
If the exponents are not all equal, the determinant is positive.
The negative Rayleigh value supplies a negative eigenvalue, so the
positive determinant forces exactly two negative eigenvalues.
If all exponents are equal, $L=-vv^T$, with one negative and two zero
eigenvalues. In both cases there is a two-dimensional subspace on
which $L$ is nonpositive.

Since $a>3$, the identity
\[
K(3)=L-\operatorname{diag}_{t\in\{a,b,c\}}
\left(3+\frac{9P}{t^2(t-3)}\right)
\]
subtracts a positive definite matrix. Thus $K(3)$ is strictly negative
on that subspace and has at least two negative eigenvalues. As $D-3I$
is positive definite, $\mathcal C-3I$ also has at least two negative
eigenvalues. The connected complement quotient has a simple zero
eigenvalue and all other eigenvalues are positive. At least one of
them therefore lies strictly between zero and three. This argument
also covers zero values among the $k_t(3)$.
\end{proof}

\section{The endpoint-one sum bound and spectral separation}
On endpoint one the sum of the two larger exponents is quadratic in the
minimum. We retain the full spectral statement and its written proof;
the sum bound and the symmetric quotient are the inputs used in the
largest-root completion.

\begin{theorem}
\label{thm:endpoint-one-spectrum}
Let $2\leq a\leq b\leq c$ be real and put $s=a+b+c$. The ordered eigenvalues of the
six-support complement quotient, counted with multiplicity, satisfy
\[
 0=\lambda_1<\lambda_2\leq\lambda_3<a+b<s
 \leq\lambda_4\leq\lambda_5\leq\lambda_6.
\]
Equality $\lambda_4=s$ holds exactly at $(a,b,c)=(2,2,2)$.
If also $h_C(1)=0$, then
\[
 b+c\geq2a^2-2a+1,
\]
with equality exactly when $b=a$ and $c=(a-1)(2a-1)$.
If in addition $a>2+\sqrt3$, root one is
simple and is the smallest positive quotient root; all four other
quotient roots are strictly greater than $a$. Their ordered roots satisfy
\[
 a<\mu_1<\min\{c,a+b\},\qquad s<\mu_2\leq\mu_3\leq\mu_4.
\]
\end{theorem}

\begin{proof}
Let $S=W^{1/2}CW^{-1/2}$, with $W=\operatorname{diag}(a,b,c,ab,ac,bc)$.
It is symmetric positive semidefinite with a one-dimensional zero kernel.
The principal block on supports $3,12,13,23$ is
\[
 \begin{pmatrix}ab+a+b&-\sqrt{abc}\\-\sqrt{abc}&c\end{pmatrix}
 \mathbin{\oplus}(b)\mathbin{\oplus}(a).
\]
The two-dimensional block has polynomial
$Q(x)=x^2-(ab+a+b+c)x+c(a+b)$, with $Q(a+b)=-ab(a+b)<0$.
Exactly three eigenvalues of this four-dimensional block lie below $a+b$;
interlacing gives $\lambda_3<a+b$.

The principal block on singleton supports $1,2$, shifted by $s$, is
\[
 \begin{pmatrix}bc-a&-\sqrt{ab}\\-\sqrt{ab}&ac-b\end{pmatrix}.
\]
Its first leading minor is $bc-a\geq a(a-1)>0$. Its determinant is
$c(abc-a^2-b^2)$, where
\[
 abc-a^2-b^2=ab(c-b)+(a-1)(b-a)(b+a)+(a-2)a^2.
\]
This is positive except at $(2,2,2)$. Away from that point,
interlacing gives $\lambda_5,\lambda_6>s$. Put $p=ab+ac+bc$, $r=abc$.
Direct expansion gives $h_C(s)=-rsV$, where
\[
 V=r(s+3)-sp=a^2(bc-b-c)+b^2(ac-a-c)+c^2(ab-a-b)>0.
\]
The three brackets are nonnegative and vanish simultaneously only
at $(2,2,2)$. Thus no eigenvalue equals $s$. At least three lie below it
and at most four. The negative determinant $s h_C(s)$ requires
an odd number below it, counting multiplicity; hence exactly three
lie below it and $\lambda_4>s$. At the boundary the characteristic
polynomial is $x(x-6)(x^2-12x+12)^2$, giving $\lambda_4=s=6$.

Now suppose $h_C(1)=0$.
Put $q=b+c$, $z=bc$ and $K=4a^2-2a+1$. The generic determinant is
\begin{align*}
 h_C(1)&=(q-K)z^2-a(a-1)(q+2a-1)z+P_a(q),\\
 P_a(q)&=(a^2-1)q^3+(a^3-a^2-3a+3)q^2
          -3(a-1)^2q-(a-1)^3.
\end{align*}
Since $(b-a)(c-a)\geq0$, we have $z\geq z_0=a(q-a)$. When
$q\leq Q:=2a^2-2a+1$, the coefficient $q-K\leq-2a^2$ is negative,
and the difference between the displayed polynomial at $z$ and at
$z_0$ is
\[
 (z-z_0)\bigl[(q-K)(z+z_0)-a(a-1)(q+2a-1)\bigr]\leq0,
\]
strictly unless $z=z_0$. Set $t=q-a\geq a$. At $z_0$ the value factors as
\[
 \bigl[(a-1)(2a-1)-t\bigr]B(a,t),
\]
where $B=a^3-2a^2t^2+a^2t+a^2+3at-3a+t^2-2t+1$. For
$t=a+u$, with $u\geq0$, the identity
\begin{align*}
 -B(a,a+u)={}&(2a^2-1)u^2\\
 &+\bigl[(a-2)(4a^2+7a+9)+20\bigr]u\\
 &+(a-2)\bigl[2a^3+a(2a-1)+3\bigr]+5
\end{align*}
has strictly positive coefficients for every $a\geq2$, so $B<0$.
Thus $q<Q$ gives $h_C(1)<0$. At $q=Q$, equality
requires $z=z_0$, hence $b=a$ by ordering, and $c=(a-1)(2a-1)$.
The factorization also verifies this equality case.

Now assume also $a>2+\sqrt3$.
The pair-support principal block is $\operatorname{diag}(c,b,a)$.
The sum bound implies
$b+c-(a^2+2a)\geq a^2-4a+1>0$, so
\[
 h_C(a)=-a^2b^2c^2(a^2+2a-b-c)>0.
\]
No eigenvalue equals $a$. The positive determinant
$\det(aI-S)=a h_C(a)$ forces an even number of eigenvalues below
$a$, counted with multiplicity. There are at least two, zero and one,
and at most three by interlacing. There are therefore exactly two:
zero and a simple one. All four remaining roots exceed $a$.

The pair-support block gives $\lambda_3\leq c$. The identity
$h_C(c)=-a^2b^2c^2(c^2+2c-a-b)<0$ excludes equality.
Together with the global total-sum gap, this proves the residual bounds.

The real cutoff is sharp. On the ordered endpoint-one family
$b=a$, $c=(a-1)(2a-1)$, we have
$h_C(a)=a^4(a-1)^2(2a-1)^2(a^2-4a+1)$.
At $a=2+\sqrt3$ a residual root equals $a$. For $2\leq a<2+\sqrt3$,
the determinant at $a$ is negative, so the same interlacing and parity
argument gives three roots below $a$; a residual root lies below $a$.
\end{proof}

\section{The endpoint-two alternative}
For integer exponents, every noninteger quotient root $\mu$ transfers to
the graph eigenvalue $N+abc-1-\mu$: the exception $N-\mu=0$ in
Lemma~\ref{lem:join} is impossible because $N$ is an integer. Conversely,
the earlier repeated-entry and small-minimum proofs exhibit noninteger
roots in the six-support space itself. The repeated-entry invariant blocks
lie in that space, and~\eqref{eq:complement} reflects their roots from $B$
to $C$. These are the quotient obstructions used for the smaller cases
in the endpoint-two theorem.

The endpoint-two surface admits a second noninteger root at large minima.
Two geometric bounds first confine the surface to a rectangle on which
positive polynomial identities give that root. A complete finite base
then handles the smaller fully distinct triples.

\begin{lemma}
\label{lem:endpoint-two-rectangle}
For ordered real exponents $4\leq a\leq b\leq c$, the equation
$h_C(2)=0$ requires $b<2a-2$. If $a\geq8$, it also requires
$c<9a/4-8$.
\end{lemma}

\begin{proof}
First put $b=2a-2+u$ and $c=b+v$, where $u,v\geq0$. Define
\begin{align*}
 A_0(u)&=2u^2+7(a-2)u+2(3a^2-14a+12),\\
 B_0(u)&=(a+2)u^3+(4a^2+10a-12)u^2\\
 &\quad +(4a^3+25a^2-48a+12)u+2(13a^3-28a^2+8a+8),\\
 D_0(u)&=2(a-2)(9a^3+24a^2-56a+16)\\
 &\quad +(33a^3+20a^2-208a+136)u
 +(20a^2+23a-62)u^2+4(a+2)u^3,\\
 E_0(u)&=2(3a^3-a^2-8a+4)+(5a^2+3a-10)u+(a+2)u^2.
\end{align*}
Direct expansion of the quotient determinant gives
\[
 h_C(2)=A_0B_0+\tfrac12(A_0B_0)'v+D_0v^2+E_0v^3,
\]
where the derivative is in $u$. Substitution $a=4+m$ makes every
coefficient of $A_0,B_0,D_0,E_0$ strictly positive for $m\geq0$.
In particular, $3a^2-14a+12=3m^2+10m+4$.
Thus $h_C(2)>0$ on $b\geq2a-2$, proving the first bound.

For $a\geq8$ and $a\leq b<2a-2$, parameterize the proposed maximum
tail by
\[
 a=8+m,\qquad b=\frac{a+(2a-2)t}{1+t},\qquad
 c=\frac{9a}{4}-8+v,\qquad m,t,v\geq0.
\]
The inverse $t=(b-a)/(2a-2-b)$ has positive denominator. The exact
identity is
\begin{align*}
 64(1+t)^3h_C(2)
 &=9m^6(1+2t)\bigl[14(1-t)^2+5t+t^2\bigr]\\
 &\quad+308m^5t(t-1)^2+P(m,t,v).
\end{align*}
The first two terms are nonnegative. The complete polynomial $P$,
listed in Appendix~\ref{sec:endpoint-two-identities}, has $84$ strictly
positive coefficients and constant $3\,014\,656$. Hence the right side
is positive, and so is $h_C(2)$. Together with the first bound, this
proves $c<9a/4-8$ on the zero surface.
\end{proof}

\begin{theorem}
\label{thm:endpoint-two-completion}
For ordered real $40\leq a\leq b\leq c$ with $h_C(2)=0$,
\[
 h_C(4)>0>h_C(5).
\]
Consequently the quotient has a noninteger root in $(4,5)$.
For every positive integer exponent triple with $h_C(2)=0$, the
quotient spectrum is nonintegral. The integer conclusion uses the
complete finite base specified below and the earlier small-minimum
certificates.
\end{theorem}

\begin{proof}
Write $s=a+b+c$, $p=ab+ac+bc$ and $r=abc$, and set
\[
 D_4=h_C(4)-3h_C(2),\qquad D_5=4h_C(2)-h_C(5).
\]
The full identities in Appendix~\ref{sec:endpoint-two-identities} give
the following positive expansions. For $a=8+m,b=8+t,c=8+u$,
$D_5$ has $44$ strictly positive coefficients and constant
$4\,629\,843$. Thus $D_5>0$ for arbitrary $a,b,c\geq8$.

Lemma~\ref{lem:endpoint-two-rectangle} confines an endpoint-two zero
at $a\geq40$ to $a\leq b<2a$ and $a\leq c<9a/4$. Put
\[
 a=40+m,\qquad b=\frac{a(1+2t)}{1+t},\qquad
 c=\frac{a(4+9u)}{4(1+u)},\qquad m,t,u\geq0.
\]
The inverse substitutions are $t=(b-a)/(2a-b)$ and
$u=4(c-a)/(9a-4c)$, with positive denominators. The polynomial
$64(1+t)^3(1+u)^3D_4$ has $112$ strictly positive coefficients and
constant $611\,305\,472$, all displayed in the appendix. Therefore
$D_4>0$. At $h_C(2)=0$ these two signs give $h_C(4)>0>h_C(5)$;
the intermediate value theorem supplies the root in $(4,5)$.

For positive integer triples, repeated exponents are nonintegral by
Theorem~\ref{thm:repeated}, and minima at most seven by
Corollary~\ref{cor:minimum-seven}. The only remaining base is
\[
 8\leq a\leq39,\qquad a<b<2a-2,\qquad b<c<\frac{9a}{4}-8.
\]
The largest permitted integer $c$ is $\lfloor(9a-33)/4\rfloor$.
This domain contains exactly $8\,658$ triples; its per-minimum counts
are listed in Appendix~\ref{sec:endpoint-two-identities}. Every endpoint
value is nonzero, with smallest absolute value $8\,208$.
The complete
\href{https://github.com/the-omega-institute/ideal-intersection-laplacian/blob/4afe151e1dc7193e6776cc1b08d3606aa804238e/results/endpoint-two-completion-base.csv}{integer certificate}
records $h_C(2)$ by integer Horner evaluation and
$\det(2I-C)=2h_C(2)$ by an independent fraction-free $6\times6$
Bareiss determinant for every triple. Thus no fully distinct endpoint-two
zero occurs in this necessary base. The real tail covers all larger
minima, completing the integer conclusion. The proof does not require
emptiness of the unbounded integer zero surface.
\end{proof}

\section{Completion of the three-prime classification}
\label{sec:three-prime-classification}

The largest quotient root supplies the final obstruction. Its open unit
interval completes endpoint one without deciding the quartic's smallest
root or enumerating integer exponent points on the endpoint surface.

\begin{theorem}
\label{thm:largest-root-unit-interval}
Let $2\leq a\leq b\leq c$ be real, with $b-a\geq3$ and $c\geq a^2-a$.
Put $s=a+b+c$ and $M=bc+s$. Then
\[
 M<\lambda_{\max}(C)<M+1.
\]
For integer exponents in this domain, the ideal intersection graph is
Laplacian nonintegral. The real interval has a written proof with no
finite exponent base.
\end{theorem}

\begin{proof}
The quotient determinant gives
\[
 h_C(M)=-a^2bc\,T(a,b,c),
\]
where
\[
 T=(2b^2-ab+b-a)c^2-[(a-1)b^2+a^2]c-ab(a+b).
\]
The quadratic coefficient of $T$ in $c$ is positive, and
\[
 T(a,b,b)=2b[(b-a+1)b^2-ab-a^2]\geq2b^3(b-a-1)>0,
 \qquad \partial_cT(a,b,b)\geq b^3>0.
\]
The derivative bound follows from
$\partial_cT(a,b,b)=b^2(4b-3a+3)-2ab-a^2$ and $a\leq b$.
Thus $T>0$ for $c\geq b$. Put $P(a,b,c)=h_C(M+1)$. For $a\geq3$, under
$a=3+m$, $b=a+3+u$, $c=a^2-a+v$, with $m,u,v\geq0$,
the complete identities in Appendix~\ref{sec:largest-root-identities}
express $T$ and $P$ as polynomials with respectively $36$ and $170$
strictly positive coefficients, including constants $1404$ and $513220$.
For $2\leq a\leq3$, write $a=2+m$, $b=a+3+u$, $c=b+v$,
where $0\leq m\leq1$ and $u,v\geq0$. The same appendix gives
$P\geq B(u,v)>0$, with all $35$ coefficients of $B$ displayed.
Thus $h_C(M)<0<h_C(M+1)$, giving a root in $(M,M+1)$.

To identify it as the largest root, use the symmetric similarity $S$
from Theorem~\ref{thm:endpoint-one-spectrum}. For $x>c$, the pair-support
block of $xI-S$ is positive definite. Its Schur complement is
\[
 H(x)=\operatorname{diag}(\ell_a,\ell_b,\ell_c)+zz^T,\qquad
 z=(\sqrt a,\sqrt b,\sqrt c)^T,\qquad
 \ell_t=x-s-\frac{xabc}{t(x-t)}.
\]
At $x=M+1$, both $x-b$ and $x-c$ exceed $bc$. Hence
\begin{align*}
 \ell_b&=c(b-a)+1-\frac{abc}{x-b}>c(b-a)+1-a>0,\\
 \ell_c&=b(c-a)+1-\frac{abc}{x-c}>b(c-a)+1-a>0.
\end{align*}
Here $c\geq b\geq a+3$. The $b,c$ principal block of $H$ is
positive definite, and $\det(xI-S)=xP>0$ gives $\det H>0$.
Its remaining scalar Schur complement is therefore positive.
Consequently $xI-S$ is positive definite, and every eigenvalue is less
than $M+1$. Since $\det(MI-S)=Mh_C(M)<0$, not all eigenvalues can
lie below $M$, so the largest lies in the stated interval. For integer exponents,
the established complement and universal-class lifts transfer its
nonintegrality to the ideal intersection graph.
\end{proof}

The first two middle gaps use exponent-root brackets rather than the
largest-root interval. The following proof records their signs explicitly
so that the classification does not rely on an unstated uniqueness claim.

\begin{lemma}
\label{lem:endpoint-one-small-gaps}
Let $a\geq2$ be real, $d\in\{1,2\}$ and $b=a+d$. Put
$L_1=2a^2-a-2$ and $L_2=2a^2+a-6$. As a polynomial in $c$,
$h_C(1)$ has exactly one root above $b$, lying in $(L_d,L_d+1)$.
Consequently, for integer $a\geq2$, no integer $c>b$ satisfies $h_C(1)=0$.
\end{lemma}
\begin{proof}
For fixed $a,b$, write $Q(c)=h_C(1)=A c^3+B c^2+D c+E$, where
\begin{align*}
 A&=a^2+b^2-1,\\
 B&=a^3-4a^2b^2+2a^2b-a^2+2ab^2+ab-3a+b^3-b^2-3b+3,\\
 D&=(a-1)(b-1)(2ab+3a+3b-3),\\
 E&=(a-1)(b-1)(a+b-1)(ab+a+b-1).
\end{align*}
Both $A$ and $E$ are positive. Put $L_1=2a^2-a-2$ and
$L_2=2a^2+a-6$. Direct expansion gives
\begin{align*}
 -Q(L_1)&=4a(a+1)(a^4-4a^2-2a+6),\\
 Q(L_1+1)&=4a(a-1)(a+1)(a^3-a^2+1),\\
 -Q(L_2)&=16a^5+72a^4-76a^3-433a^2+69a+595,\\
 Q(L_2+1)&=4(a+2)(2a^5-2a^4-17a^3+25a^2+28a-42),
\end{align*}
where each pair uses its own $b=a+d$. The complete positive vectors
at $a=2+m$, including those for $-Q(b)$, are in
Appendix~\ref{sec:largest-root-identities}. Thus $Q(0)>0>Q(b)$ and
$Q(L_d)<0<Q(L_d+1)$. Moreover,
\[
 L_1-b=1+6(a-2)+2(a-2)^2>0,\qquad
 L_2-b=2(a-2)(a+2)\geq0.
\]
The second equality occurs only at $a=2$; the final open interval
still lies strictly above $b$ there.
The cubic has one root in each of $(-\infty,0)$, $(0,b)$ and
$(L_d,L_d+1)$. These three distinct roots exhaust its degree.
The unique root above $b$ is therefore in $(L_d,L_d+1)$.
For integer $a$, the endpoints are consecutive integers, proving the
nonvanishing conclusion without a finite base.
\end{proof}

\begin{theorem}[Three-prime nonintegrality]
\label{thm:three-prime-classification}
For distinct primes $p,q,r$ and positive integers $a,b,c$, the ideal
intersection graph of $\mathbb Z_{p^a q^b r^c}$ is Laplacian nonintegral.
This global conclusion combines written proofs with the established
small-minimum and endpoint-two finite certificates.
\end{theorem}
\begin{proof}
Sort the exponents. Repeated entries are covered by
Theorem~\ref{thm:repeated}; minima at most seven by
Corollary~\ref{cor:minimum-seven}. Suppose $8\leq a<b<c$ and the
spectrum is integral. Theorem~\ref{thm:uniform-low-root} supplies a
positive quotient root in $(0,3)$, necessarily one or two.
Theorem~\ref{thm:endpoint-two-completion} excludes the second alternative,
so $h_C(1)=0$.

For $b-a=1$ or $2$, Lemma~\ref{lem:endpoint-one-small-gaps} contradicts
integer $c>b$. Otherwise $b-a\geq3$, and
Theorem~\ref{thm:endpoint-one-spectrum} gives
$b+c\geq2a^2-2a+1$. Since $c\geq b$, we have $c>a^2-a$.
Theorem~\ref{thm:largest-root-unit-interval} supplies a noninteger
quotient root, again contradicting integrality. These cases exhaust
all positive exponent triples.

The minimum-through-seven theorem retains its $27\,562$-triple base
and earlier small-minimum dependencies; the endpoint-two theorem
retains its $8\,658$-triple base. The new largest-root and small-gap
steps have written unbounded proofs. No additional exponent scan or
integer-point enumeration is used in this completion.
\end{proof}
